% authority: science_draft
\documentclass[11pt]{article}

\usepackage[margin=1in]{geometry}
\usepackage{amsmath,amssymb,amsthm}
\usepackage{booktabs}
\usepackage{longtable}
\usepackage{enumitem}

\newcommand{\EqSrc}{\mathrm{EqSrc}}
\newcommand{\EqSrcT}{\mathrm{EqSrc}_{T}}
\newcommand{\Block}{\mathrm{Block}}
\newcommand{\NotEqSrc}{\mathrm{not\_EqSrc}}
\newcommand{\RetainH}{\mathrm{RetainH}}
\newcommand{\GenH}{\mathrm{GenH}}
\newcommand{\Fsrc}{\mathcal{F}_{\mathrm{src}}}
\newcommand{\Osrc}{\mathcal{O}_{\mathrm{src}}}
\newcommand{\Msrc}{\mathcal{M}_{\mathrm{src}}}
\newcommand{\Isrc}{\mathcal{I}_{\mathrm{src}}}
\newcommand{\Asrc}{\mathcal{A}_{\mathrm{src}}}
\newcommand{\Psrc}{\mathcal{P}_{\mathrm{src}}}
\newcommand{\Nsrc}{\mathcal{N}_{\mathrm{src}}}
\newcommand{\Csrc}{C_{\mathrm{src}}}
\newcommand{\NoTarget}{\mathsf{NoTarget}}

\newtheorem{definition}{Definition}
\newtheorem{theorem}{Theorem Candidate}
\newtheorem{countermodel}{Minimal Countermodel Slot}
\newtheorem{remark}{Boundary Remark}

\title{EqSrc Family-Closure Theorem or Countermodel v1}
\author{Ontology Formalizer Packet RT-20260707-020}
\date{2026-07-07}

\begin{document}
\maketitle

\section{Control Status}

This artifact implements v18 P3-T02 as a draft/control science artifact. It
executes one bounded EqSrc family-closure theorem-or-countermodel attempt over
the typed source family introduced by
\(\mathrm{SourceEquivalenceTypedObject\_v1}\).

The primary result is:
\[
  \texttt{family\_closure\_theorem\_candidate\_supplied}.
\]

\begin{verbatim}
primary_result: family_closure_theorem_candidate_supplied
\end{verbatim}

The required minimal countermodel slot attempted is:
\[
  \texttt{missing\_inverse\_countermodel}.
\]

The result has no proof-promotion authority. It does not discharge general
\(\EqSrc\), adopt \(\RetainH\), adopt \(\GenH\), adopt a source law, edit
canonical ontology, adopt \(\mathrm{MetricData}(E)\), expand
\(g_{\mathrm{eff}}\), derive matter coupling, derive Einstein equations,
promote a benchmark, issue a Gate Chair verdict, authorize external outreach,
or claim completed derivation.

\section{Prior Record-Local \(\EqSrc\) Theorem}

The prior record-local theorem candidate works inside one declared finite
source-record orbit. It supplies record-local relabeling witnesses and checks
identity, inverse, and composition inside that local orbit. It does not by
itself establish a theorem over an arbitrary source family.

The present packet treats the record-local theorem as an input pattern. It asks
which source-side hypotheses are sufficient to lift the equivalence-relation
argument to the typed family \(\Fsrc\), and which minimal countermodel appears
when one hypothesis is removed.

\section{Typed Source Family}

Let
\[
  \Fsrc=(\Osrc,\Msrc,\Isrc,\Csrc,\Asrc,\Psrc,\Nsrc)
\]
be a declared source-side family, where:
\begin{enumerate}[label=(\alph*)]
  \item \(\Osrc\) is a class or bounded family fragment of source objects;
  \item \(\Msrc\) is a class of admissible source morphisms or relabelings;
  \item \(\Isrc\) is a source invariant ledger;
  \item \(\Csrc:\Osrc\times\Osrc\to\{\EqSrc,\NotEqSrc,\Block\}\) is a
  source-only comparison rule;
  \item \(\Asrc\) records source paths, object ids, hashes, scopes, and
  dependencies;
  \item \(\Psrc\) records forbidden-channel and proxy-edge incidence;
  \item \(\Nsrc\) is a negative-control set.
\end{enumerate}

Every component is source-side. None is a target topology, target metric,
detector protocol, empirical calibration, stress-energy object, matter action,
benchmark behavior, validator status, registry authority, role identity,
handoff state, commit state, or generated derivative.

\section{Family Invariant Ledger}

For a morphism \(f:X\to Y\) in \(\Msrc\), write
\[
  f \models \mathrm{Pres}(\Isrc,\Asrc,\Psrc)
\]
when \(f\) preserves the relevant invariant entries, annotation entries, and
forbidden-channel or proxy-edge incidence by an explicit source-side rule.

The family-level theorem candidate uses the following ledger hypothesis:

\[
  \forall f\in\Msrc,\quad
  f \models \mathrm{Pres}(\Isrc,\Asrc,\Psrc)
  \text{ or } \Csrc(\mathrm{dom}(f),\mathrm{cod}(f))=\Block.
\]

This is a hypothesis, not a completed ledger proof. The current project record
does not promote it to a source law.

\section{Candidate Family-Level \(\EqSrc\) Relation}

\begin{definition}[Typed family relation]
For \(X,Y\in\Osrc\), define
\[
  X\,\EqSrcT\,Y
\]
when there exists a witness \(w_{XY}:X\to Y\) in \(\Msrc\) such that:
\begin{enumerate}[label=(\roman*)]
  \item \(w_{XY}\models \mathrm{Pres}(\Isrc,\Asrc,\Psrc)\);
  \item \(\Csrc(X,Y)=\EqSrc\);
  \item all negative-control checks in \(\Nsrc\) remain source-side and
  fail-closed;
  \item \(\NoTarget(w_{XY})\) holds.
\end{enumerate}
\end{definition}

\(\NoTarget\) means that the witness and comparison do not use target
topology, target metric, proper time, detector protocol, stress-energy,
matter action, Einstein equations, benchmark behavior, validator status,
registry metadata, role identity, handoff state, commit state, or generated
derivatives as premises.

\section{Identity Closure Attempt}

The identity hypothesis is:
\[
  \forall X\in\Osrc,\quad \mathrm{id}_{X}\in\Msrc
  \quad\text{and}\quad
  \mathrm{id}_{X}\models \mathrm{Pres}(\Isrc,\Asrc,\Psrc).
\]

If also \(\Csrc(X,X)=\EqSrc\) under the no-target guard, then
\[
  X\,\EqSrcT\,X.
\]

Thus reflexivity follows from supplied source-side identity witnesses. It is
not inferred from process authority or from target-side behavior.

\section{Inverse Closure Attempt}

The inverse hypothesis is:
\[
  \forall w_{XY}:X\to Y\in\Msrc,\quad
  \exists w_{YX}:Y\to X\in\Msrc
\]
such that \(w_{YX}\models \mathrm{Pres}(\Isrc,\Asrc,\Psrc)\),
\(\Csrc(Y,X)=\EqSrc\), and \(\NoTarget(w_{YX})\) holds whenever
\(w_{XY}\) witnesses \(X\,\EqSrcT\,Y\).

If the inverse witness is missing, symmetry is not derivable. That missing
case is the minimal countermodel slot recorded below.

\section{Composition Closure Attempt}

The composition hypothesis is:
\[
  w_{XY}:X\to Y,\quad w_{YZ}:Y\to Z\in\Msrc
  \quad\Longrightarrow\quad
  w_{YZ}\circ w_{XY}:X\to Z\in\Msrc,
\]
with source-side preservation of \(\Isrc,\Asrc,\Psrc\), \(\Csrc(X,Z)=\EqSrc\),
and \(\NoTarget(w_{YZ}\circ w_{XY})\).

If composition is not supplied or derived, transitivity is not derivable.

\section{\(\RetainH\) Trigger Analysis}

\(\RetainH\) is not required for the conditional theorem over an already
declared closed family \(\Fsrc\). The theorem candidate asks only whether
identity, inverse, composition, ledger stability, comparison, and no-target
guards are present inside the declared family.

\(\RetainH\) becomes a live boundary condition when a later packet asks
whether source-record structure, certificate status, comparison status,
negative controls, or invariant status is preserved under an \(H\)-indexed
retention operation. The current sources do not derive that primitive.

Classification:

\begin{verbatim}
RetainH_status_for_this_theorem: not_required_here
RetainH_status_for_H_retention_extension: candidate_definition_needed
RetainH_adopted: false
adoption_requested: false
\end{verbatim}

\section{\(\GenH\) Trigger Analysis}

\(\GenH\) is not required for the conditional theorem over an already declared
closed family \(\Fsrc\). The theorem candidate does not claim to generate all
family members; it assumes the family and its closure witnesses are declared.

\(\GenH\) becomes a live boundary condition when a later packet constructs,
enumerates, or quantifies over an \(H\)-generated source family. The current
sources do not derive that primitive.

Classification:

\begin{verbatim}
GenH_status_for_this_theorem: not_required_here
GenH_status_for_H_generated_extension: candidate_definition_needed
GenH_adopted: false
adoption_requested: false
\end{verbatim}

\section{Minimal Countermodel Search}

\begin{countermodel}[Missing inverse countermodel]
Let
\[
  \Osrc=\{A,B\},\qquad
  \Msrc=\{\mathrm{id}_{A},\mathrm{id}_{B},f:A\to B\}.
\]
Assume all listed morphisms are source-side and preserve the chosen invariant
and annotation entries. Define the comparison rule by
\[
  \Csrc(A,A)=\EqSrc,\quad
  \Csrc(B,B)=\EqSrc,\quad
  \Csrc(A,B)=\EqSrc,\quad
  \Csrc(B,A)=\Block.
\]
The block in the last entry is caused only by the absence of a source-side
inverse \(f^{-1}:B\to A\) in \(\Msrc\).
\end{countermodel}

This is a finite source-only witness. It does not import target topology,
target metric, detector protocol, stress-energy, matter action, Einstein
equations, benchmark behavior, or process authority. It shows that inverse
closure is a real theorem hypothesis: \(A\,\EqSrcT\,B\) is witnessed by \(f\),
but \(B\,\EqSrcT\,A\) is blocked. Therefore symmetry fails without the inverse
hypothesis.

Other countermodel slots remain available for later stress:
\[
\texttt{missing\_composition\_countermodel},\quad
\texttt{invariant\_ledger\_not\_family\_stable\_countermodel},\quad
\texttt{RetainH\_needed\_countermodel},\quad
\texttt{GenH\_needed\_countermodel}.
\]

\section{Theorem Candidate or Obstruction}

\begin{theorem}[Conditional family closure for \(\EqSrcT\)]
Assume a declared source family \(\Fsrc\) satisfies all of the following:
\begin{enumerate}[label=(H\arabic*)]
  \item source-only object domain: every \(X\in\Osrc\) is a source object;
  \item identity closure: \(\mathrm{id}_{X}\in\Msrc\) for every \(X\);
  \item inverse closure: every witness \(w_{XY}\) has a source-side inverse
  witness \(w_{YX}\);
  \item composition closure: composable witnesses have source-side
  compositions in \(\Msrc\);
  \item invariant-ledger stability: witnesses and their inverses and
  compositions preserve \(\Isrc,\Asrc,\Psrc\);
  \item source-only comparison: \(\Csrc\) uses only source data and returns
  \(\Block\) on missing ledger or proxy-incidence evidence;
  \item no-target guard: \(\NoTarget\) holds for all witnesses and comparison
  premises.
\end{enumerate}
Then \(\EqSrcT\) is reflexive, symmetric, and transitive on \(\Osrc\).
\end{theorem}

\begin{proof}
For reflexivity, fix \(X\in\Osrc\). By H2, \(\mathrm{id}_{X}\in\Msrc\).
By H5 and H7, the identity witness preserves the source ledger and uses no
forbidden target or process-authority premise. By H6, \(\Csrc(X,X)=\EqSrc\)
when those source-side checks pass. Hence \(X\,\EqSrcT\,X\).

For symmetry, assume \(X\,\EqSrcT\,Y\), witnessed by
\(w_{XY}:X\to Y\). By H3, there is a source-side inverse witness
\(w_{YX}:Y\to X\). By H5, H6, and H7, this inverse preserves the ledger,
passes the same source-only comparison rule, and uses no forbidden target or
process-authority premise. Hence \(Y\,\EqSrcT\,X\).

For transitivity, assume \(X\,\EqSrcT\,Y\) and \(Y\,\EqSrcT\,Z\), witnessed by
\(w_{XY}\) and \(w_{YZ}\). By H4, \(w_{YZ}\circ w_{XY}\in\Msrc\). By H5, H6,
and H7, the composite preserves the source ledger, passes the source-only
comparison rule, and uses no forbidden target or process-authority premise.
Hence \(X\,\EqSrcT\,Z\).

The proof is conditional on H1--H7. It does not show that the current ontology
derives those hypotheses for every future source family.
\end{proof}

The obstruction is therefore sharpened rather than globally discharged:
`OB-P6T02-GENERAL-EQSRC-FAMILY-CLOSURE-MISSING` is reduced to the need for
source-side proof or construction of H1--H7 in each intended family, plus
separate extraction of `RetainH` and `GenH` boundaries for \(H\)-retained or
\(H\)-generated extensions.

\section{Distance-to-GR Effect}

\begin{longtable}{p{0.30\linewidth}p{0.56\linewidth}}
\toprule
Burden & Status after this packet \\
\midrule
Source ontology primitives & Unchanged; no canonical ontology edit. \\
Source equivalence EqSrc & Conditional theorem candidate supplied; general family closure not adopted and no ledger delta. \\
RetainH & Not adopted; boundary extracted as candidate-definition-needed for H-retention extensions. \\
GenH & Not adopted; boundary extracted as candidate-definition-needed for H-generated family extensions. \\
ObsLoc\_lc & Unchanged. \\
Resp\_lc & Unchanged. \\
M\_src & Unchanged. \\
g\_eff & Unchanged; no MetricData(E) or metric-scope expansion. \\
matter coupling & Blocked and unchanged; no coupling derivation or adoption. \\
Einstein equations & Blocked and unchanged; no field-equation derivation. \\
finite-variation robustness & Unchanged; no global robustness theorem. \\
benchmark promotion & Blocked; no benchmark evidence or promotion. \\
Gate Chair status & Unchanged; no Gate Chair verdict. \\
current route freeze or hard-fail status & Not frozen; route continues to P3-T03 extraction, then audit and stress. \\
\bottomrule
\end{longtable}

\begin{verbatim}
distance_to_gr_delta:
  changed: false
  effect: no_distance_delta
  ledger_row_updated: false
  reason: P3-T02 supplies a conditional theorem candidate and one minimal
    inverse-countermodel slot only. It does not prove the closure hypotheses
    from current ontology and does not adopt RetainH GenH or EqSrc.
\end{verbatim}

\section{Forbidden Conclusions}

This artifact does not authorize:

\begin{itemize}[leftmargin=*]
  \item canonical ontology edit;
  \item general \(\EqSrc\) discharge;
  \item \(\RetainH\) adoption;
  \item \(\GenH\) adoption;
  \item source-law adoption;
  \item source detector/readout semantics adoption;
  \item \(\mathrm{MetricData}(E)\) adoption;
  \item \(g_{\mathrm{eff}}\) adoption or scope expansion;
  \item matter-coupling derivation or adoption;
  \item stress-energy semantics, stress-energy tensor, or matter action;
  \item Einstein-equation derivation;
  \item benchmark promotion;
  \item Gate Chair verdict;
  \item external outreach;
  \item completed derivation;
  \item future source-extension impossibility;
  \item global no-go conclusion.
\end{itemize}

\section{Source Materials}

The AEther-Flow Research Project. (2026a). \emph{Recommendations
implementation plan continue task v18} [Internal implementation plan].
\texttt{implementations\_plans/recommendations\_implementation\_plan\_continue\_task-v18.md}.

The AEther-Flow Research Project. (2026b). \emph{Source-equivalence typed
object v1} [Research-control TeX artifact].
\texttt{research\_control/tasks/RT-20260707-015/artifacts/source\_equivalence\_typed\_object\_v1.tex}.

The AEther-Flow Research Project. (2026c). \emph{V18 P2-T05
source-equivalence typed object Refuter stress v1} [Research-control TeX
artifact].
\texttt{research\_control/tasks/RT-20260707-017/artifacts/source\_equivalence\_typed\_object\_refuter\_stress\_v1.tex}.

The AEther-Flow Research Project. (2026d). \emph{EqSrc family-closure packet
setup v1} [Research-control artifact].
\texttt{research\_control/tasks/RT-20260707-019/artifacts/eqsrc\_family\_closure\_packet\_setup\_v1.md}.

\end{document}
